\section{Methodology}
\label{sec:methodology}

\subsection{Baseline Framework}

We build on the LTR scheduling framework of~\cite{fu2024efficient}, which
trains an OPT-125M predictor with a ListMLE ranking loss to score prompts by
their expected relative output length, and uses that score to reorder the
admission queue before requests are submitted to vLLM's generation engine.
This framework operates entirely outside vLLM's own continuous-batching
scheduler: the ranking step pre-computes an admission order for a batch of
requests, and vLLM's engine handles the resulting iteration-level batching and
memory management unmodified, as shown in Figure~\ref{fig:architecture_pipeline}.

\subsection{Proposed Extension: Composite Scoring}

Table~\ref{tab:scheduling_policies} summarizes the four scheduling policies
compared in this work. For the proposed extension, each request's admission
priority is computed as
\begin{equation}
    \mathrm{score}(x) = \alpha \, s_{\mathrm{LTR}}(x) \;-\; \beta \, \frac{\mathrm{SCALE}}{\mathrm{len}(x)}
    \label{eq:composite-score}
\end{equation}
where $s_{\mathrm{LTR}}(x)$ is the OPT-125M predictor's learned score for
prompt $x$, $\mathrm{len}(x)$ is the number of tokens in $x$ under the serving
model's tokenizer, $\mathrm{SCALE}$ is a fixed constant that brings the
length term onto a comparable numeric range to $s_{\mathrm{LTR}}(x)$, and
$\alpha, \beta$ are fixed weights with $\alpha + \beta = 1$. The length term
is subtracted rather than added, so that longer prompts receive a
comparatively lower composite score: this reflects the modeling assumption
that longer prompts are more likely to elicit longer, more involved
completions, which should not be prioritized ahead of shorter prompts absent
strong evidence to the contrary from $s_{\mathrm{LTR}}$. Requests are then
admitted in decreasing order of $\mathrm{score}(x)$. Because
$\mathrm{len}(x)$ is obtained directly from tokenization, this signal adds no
additional model inference cost beyond the single forward pass already
required to compute $s_{\mathrm{LTR}}(x)$.

The predictor used to compute $s_{\mathrm{LTR}}(x)$ for the extension is
retrained with a pairwise margin ranking loss rather than the original
ListMLE loss, so that the underlying score distribution is well suited to
combination with the length term; both training configurations share the same
OPT-125M backbone and training set. Table~\ref{tab:hyperparameters} lists the
composite-scoring weights and predictor training hyperparameters used
throughout the evaluation in Section~\ref{sec:evaluation}.

\input{tables/scheduling_policies}
\input{tables/hyperparameters}

\subsection{Implementation Correction}
\label{subsec:sign-correction}

An early implementation of Eq.~\eqref{eq:composite-score} added the length
term rather than subtracting it, which inverted its intended effect and
caused longer prompts to receive a higher composite score. This was corrected
prior to the experiments reported in Section~\ref{sec:evaluation} by changing
the term's sign to match Eq.~\eqref{eq:composite-score}; we note the
correction here for transparency, since it affects how the composite score
should be interpreted relative to any earlier, uncorrected version of this
codebase.
